
\exercise
Suppose $T$ is a function from $V$ to $W\3$. The \emph{graph} of $T$\index{graph of a linear map} is the subset of
$V\1 \times W$ defined by
\[
\text{graph of}\ T = \br[\big]{ (v, Tv) \in V\1 \times W: v \in V }.
\]
Prove that $T$ is a linear map if and only if the graph of $T$
is a subspace of $V\1 \times W\3$.
\exercisecomment{Formally, a function\/ $T$ from\/ $V$ to\/ $W$ is a subset\/ $T$ of\/ $V\1 \times W$ such that for each\/ $v \in V\1$, there exists exactly one element\/ $(v, w) \in T\3$. In other words, formally a function is what is called above its graph. We do not usually think of functions in this formal manner. However, if we do become formal, then this exercise could be rephrased as follows: Prove that a function\/ $T$ from\/ $V$ to\/ $W$ is a linear map if and only if\/ $T$ is a subspace of\/ $V\1 \times W\3$.}


\exercise
Suppose that $V\1_1, \dots, V\1_m$ are vector spaces such that $V\1_1 \times \dots \times V\1_m$ is finite-dimensional. Prove that $V\1_k$ is finite-dimensional for each $k = 1, \dots, m$.


\begin{comment}
\exercise
Give an example of a vector space $V$ and subspaces $V\1_1, V\1_2$ of $V$ such that $V\1_1 \times V\1_2$ is isomorphic to $V\1_1 + V\1_2$ but $V\1_1 + V\1_2$ is not a direct sum.
\hint{The vector space $V$ must be infinite-dimensional.}

\end{comment}

\exercise
Suppose $V\1_1, \dots, V\1_m$ are vector spaces. Prove that $\L(V\1_1 \times \dots \times V\1_m, W)$ and $\L(V\1_1, W) \times \dots \times \L(V\1_m, W)$ are isomorphic vector spaces.
\exercisecomment{There is no assumption in the exercise above or in the two following exercises that the vector spaces are finite-dimensional.}



\exercise
Suppose $W\3_1, \dots, W\3_m$ are vector spaces. Prove that $\L(V\1, W\3_1 \times \dots \times W\3_m)$ and
$\L(V\1, W\3_1) \times \dots \times \L(V\1, W\3_m)$ are isomorphic vector spaces.


\exercise
For $m$ a positive integer, define $V^m$ by\sindex{Vm@$V^m\mskip -2mu$}
\[
V^m = \underbrace{V\1 \times \dots \times V}_{m \text{\ times}}\!.
\]
Prove that $V^m$ and $\L\paren[\big]{\FF{m}, V}$ are isomorphic vector spaces.


\exercise
Suppose that $v, x$ are vectors in $V$ and that $U, W$ are subspaces of $V$ such that $v + U = x + W\3$. Prove that $U = W\3$.


\exercise
Let $U = \br[\big]{ (x, y, z) \in \R3: 2x + 3y + 5z= 0}$. Suppose $A \subset \RR3$. Prove that $A$ is a translate of $U$ if and only if there exists $c \in \R{}$ such that\label{3translates}
\[
A = \br[\big]{ (x, y, z) \in \R3: 2x + 3y + 5z= c}.
\]
\exercisedisplayend


\exercise
\begin{SUBtheorem}
\item
Suppose $T \in \L(V, W)$ and $c \in W\3$. Prove that $\br{x \in V : Tx = c}$ is either the empty set or is a translate of
$\ker T\3$.

\item
Explain why the set of solutions to a system of linear equations such as \ref{3inhomeq} is either the empty set or is a translate of some subspace of $\FF n$.
\end{SUBtheorem}


\exercise
Prove that a nonempty subset $A$ of $V$ is a translate of some subspace of $V$ if and only if $\la v + (1-\la)w \in A$ for all $v, w \in A$ and all $\la \in \F{}\3$. \label{affine}


\exercise
Suppose $A_1 = v + U_1$ and $A_2 = w + U_2$ for some $v, w \in V$ and some subspaces $U_1, U_2$ of $V\1$. Prove that the intersection $A_1 \cap A_2$ is either a translate of some subspace of $V$ or is the empty set.

\pagebreak[4]


\exercise
Suppose $U = \br[\big]{(x_1, x_2, \dots) \in \F{\infty}: x_k \ne 0\ \text{for only finitely many}\ k}$.
\begin{subtheorem}
\item
Show that $U$ is a subspace of $\FF{\infty}$.
\item
Prove that $\FF{\infty}/U$ is infinite-dimensional.
\end{subtheorem}
\exercisesubtheoremend


\exercise
Suppose $v_1, \dots, v_m \in V\1$. Let
\[
A = \{ \la_1 v_1 + \dots + \la_m v_m : \la_1, \dots, \la_m \in \F{}\ \text{and}\ \la_1 + \dots + \la_m = 1 \}.
\]
\begin{subtheorem}*
\item
Prove that $A$ is a translate of some subspace of $V\1$.
\item
Prove that if $B$ is a translate of some subspace of $V$ and $\br{v_1, \dots, v_m} \subset B$, then $A \subset B$.
\item
Prove that $A$ is a translate of some subspace of $V$ of dimension less than~$m$.
\end{subtheorem}
\exercisesubtheoremend


\exercise
Suppose $U$ is a subspace of $V$ such that $V\1/U$ is finite-dimensional. Prove that $V$ is isomorphic to
$U \times (V\1/U)$.


\exercise
Suppose $U$ and $W$ are subspaces of $V$ and $V = U \oplus W\3$. Suppose $w_1, \ldots, w_m$ is a basis of $W\3$. Prove that
$w_1 + U, \ldots, w_m + U$ is a basis of $V\1/U$.


\exercise
Suppose $U$ is a subspace of $V$ and $v_1 + U, \dots, v_m + U$ is a basis of $V\1/U$ and $u_1, \dots, u_n$ is a basis of $U$. Prove that $v_1, \dots, v_m, u_1, \dots, u_n$ is a basis of $V\1$. \label{basisUVmodU}


\exercise
Suppose $\phi \in \L(V\1, \F{})$ and $\phi \ne 0$. Prove that $\dim V \!/\! (\ker \phi) = 1$.


\exercise
Suppose $U$ is a subspace of $V$ such that $\dim V\1/U = 1$. Prove that there exists $\phi \in \L(V\1, \F{})$ such that $\ker \phi = U$.


\exercise
Suppose that $U$ is a subspace of $V$ such that $V\1/U$ is finite-dimensional.
\begin{subtheorem}
\item
Show that if $W$ is a finite-dimensional subspace of $V$ and \mbox{$V = U + W$}, then $\dim W \ge \dim V\1/U$.
\item
Prove that there exists a finite-dimensional subspace $W$ of $V$ such that \mbox{$\dim W = \dim V\1/U$} and \mbox{$V = U \oplus W$}.
\end{subtheorem}
\exercisesubtheoremend


\exercise
Suppose $T \in \L(V\1, W)$ and $U$ is a subspace of $V\1$. Let $\pi$ denote the quotient map from $V$ onto $V\1/U$. Prove that there exists $S \in \L(V\1/U, W)$ such that $T = S \circ \pi$ if and only if $U \subset \ker T\3$. \label{1715}


\begin{comment} %This exercise deleted from LADR4e because it is too similar to Exercise 18 in Section 3.E of LADR3e.
\exercise
Suppose $U$ is a subspace of $V\1$. Define $\Gamma \colon \L(V\1/U, W) \to \L(V\1, W)$ by
\[ \addtolength{\belowdisplayskip}{-3bp}
\Gamma(S) = S \circ \pi.
\]
\begin{subtheorem}
\item
Show that $\Gamma$ is a linear map.
\item
Show that $\Gamma$ is injective.
\item
Show that
\mbox{$\ran \Gamma = \{ T \in \L(V\1, W): Tu = 0 \text{\ for every\ } u \in U\}$}.
\end{subtheorem}

\end{comment}

\begin{comment} % Exercise deleted because T^2 has not yet been defined.
\exercise
Suppose $T \in \L(V)$ and $\pi$ is the quotient map of $V$ onto $V\1/(\ker T)$. Prove that
\[
\ker \pi \circ \tilde T = \br{v + \ker T : v \in V \text{ and } T^2 v = 0}.
\]
\exercisecomment{The exercise above implies that\/ $\pi \circ \tilde T$ is the zero map from\/ $V\1/(\ker T)$ to itself if and only if\/ $T^2$ is the zero map from\/ $V\1$ to itself.}


\end{comment}

